\section{Transfer indices in other integer radices}\label{radix:section}

The constructive upper bounds extend to every integer radix by the same
divisible-lattice argument used in the binary case. A separate affine
record makes its hypotheses explicit. Odd radices admit a one-threshold
record because addition of representable numbers cannot produce a
rounding midpoint. The latter fact is known~\cite[Lemma~19]{roux2014}.
This appendix gives an extension of the whole-format index, rather than
an independent novelty claim for odd-radix rounding or floor identities.
The binary information lower bounds proved earlier are not extended here.

Fix an integer radix $\beta\geq2$, precision $p\geq2$, and $E\geq1$
normal magnitude intervals. In least-subnormal units, put $n=\beta^{p-1}$
and define the finite numerical format
\begin{equation}\label{radix:eq:format}
 \mathcal F_{\beta,p,E}
 =\{0\}\cup\{\pm k:1\leq k<n\}
 \cup\{\pm k\beta^j:n\leq k<\beta n,\ 0\leq j<E\}.
\end{equation}
Its size is $N=2n((\beta-1)E+1)-1$. Every input and operand belongs to
this format, and each addition is rounded separately to nearest with
ties to an even significand. Gradual underflow is included, signed zeros
are identified, and NaNs and exception flags are excluded. Both
infinities are absorbing. Write
\[
 M=(\beta n-1)\beta^{E-1},\qquad
 U=\beta n\beta^{E-1},\qquad \tau=\frac{M+U}{2}.
\]
For even $\beta$, exact sums of magnitude at least $\tau$ overflow.
For odd $\beta$, $\tau$ is half-integral and reachable sums are integral;
overflow occurs for magnitude greater than $\tau$. Its tie case is
unreachable. Odd-radix formats here are mathematical models, not IEEE
754 formats.

The format is public. Operation counts below start from its constants
$n,M,\tau$ and ordered region tables. If only $\beta,p,E$ are supplied,
these tables can first be prepared using $O(E+\log p)$ integer
multiplications: compute $n$ by exponentiation by squaring and successive
radix powers by multiplication by $\beta$. This one-time preparation
cost is separate from the stated block-construction bounds.

For odd $\beta$, adjacent finite values differ by an odd integer
$\beta^j$, so every midpoint is half-integral. This includes boundaries
of magnitude intervals and the gradual-underflow region. Consequently,
no reachable addition is a tie. If $R_h$ is uniform nearest rounding to
$h\mathbb Z$, $h=\beta^j$, then on integer arguments
\[
 R_h(z+\lambda h)=R_h(z)+\lambda h,
 \qquad
 \lambda=\begin{cases}2,&\beta\text{ even},\\1,&\beta\text{ odd}.
 \end{cases}
\]
For even radix the two-step translation preserves tie parity; for odd
radix a one-step translation suffices because there are no ties.

Split the finite input ranks into the $2E+1$ original arithmetic
progressions: the negative normal intervals, the central interval
containing subnormals and zero, and the positive normal intervals.
Retain these boundaries in every index. On a finite-output piece with
inputs $x_0+di$, $0\leq i<L$, use either an affine record
$f(i)=c+di$, or a staircase record
\begin{equation}\label{radix:eq:record}
 f(i)=c+\lambda H\left\lfloor\frac{i}{P}\right\rfloor
       +H\sum_{s=1}^{\lambda}\mergeind{i\bmod P\geq t_s},
 \qquad P=\frac{\lambda H}{d},
\end{equation}
where $d,H$ are powers of $\beta$, $H\geq d$, and
$1\leq t_1\leq\cdots\leq t_\lambda\leq\min(P,L)$.
Thresholds may coincide or be clipped to $L<P$. Infinite-output pieces
are constant. In binary, an affine record can also be written using the
fictitious height $d/2$ of \eqref{merge:eq:record}; the explicit tag avoids
using that height in radix-power divisibility arguments.

For odd radix, \eqref{radix:eq:record} is simply the shifted-floor formula
\begin{equation}\label{radix:eq:floor}
 f(i)=c+H\left\lfloor\frac{i+P-t_1}{P}\right\rfloor.
\end{equation}
Indeed, writing $i=qP+r$, $0\leq r<P$, leaves the indicator $r\geq t_1$.
Thus one threshold does not introduce a new class of elementary functions.
The nested-floor and upper-level-set identities underlying composition
are classical; compare~\cite[Lemma~1]{ref25}. More generally, normalized
staircase extensions are nondecreasing integer maps $g$ satisfying
$g(i+P)=g(i)+\lambda$. Such periodic maps and the reverse-order
composition law for their generalized inverses occur in
\cite[Definition~2.2 and Proposition~4.1]{connesconsani2013}.
The local argument below checks the additional lattice and range-coset
conditions required by the rounding representation.

\begin{lemma}[Local tagged closure]\label{radix:lem:closure}
Suppose an interval $I$ in an $A$-record has its finite image in one
$B$-record. Then $B\circ A$ on $I$ has the same tagged form, with input
spacing $d_A$. If $B$ is affine, retain the tag and height of $A$.
If $A$ is affine and $B$ is a staircase, the composite is affine when
$d_A>H_B$, and also when $d_A=H_B$ in odd radix; otherwise its height is
$H_B$. If both are staircases, its height is $\max(H_A,H_B)$.
\end{lemma}

\begin{proof}
An affine $B$ extends numerically as $b(y)=y+\text{constant}$, so it
translates $A$. For a staircase $B$, extend \eqref{radix:eq:record} to all
real numerical inputs using $u=(y-y_B)/d_B$, $q=\lfloor u/P_B\rfloor$,
and $r=u-qP_B$. This extension is nondecreasing, has range in
$c_B+H_B\mathbb Z$, and satisfies
$b(y+\lambda H_B)=b(y)+\lambda H_B$. Coinciding and clipped thresholds
preserve these properties.

If $A$ is affine and $d_A>H_B$, the ratio $d_A/H_B$ is a positive power
of $\beta$, hence divisible by $\lambda$. Each input step therefore
advances $B$ by $d_A$. Equality works in odd radix as well. In the
remaining affine case, the composite has index period $\lambda H_B/d_A$,
increment $\lambda H_B$, and range in a coset of $H_B\mathbb Z$.

If both records are staircases, extend $A$ to integer indices and put
$H=\max(H_A,H_B)$. The number $\lambda H$ is a multiple of both
numerical periods. Thus the composite $G$ satisfies
\[
 G(i+\lambda H/d_A)=G(i)+\lambda H.
\]
When $H_B\geq H_A$, its range belongs to $G(0)+H_B\mathbb Z$.
When $H_A>H_B$, the ratio $H_A/H_B$ is a positive radix power and a
multiple of $\lambda$. All differences of $A$-outputs are multiples of
$H_A$, and $B$'s translation identity preserves these differences
exactly. The range then belongs to $G(0)+H_A\mathbb Z$.

In each nonaffine case, monotonicity, the range coset, and the period
identity imply at most $\lambda$ successive level crossings per period.
Their first indices give the thresholds. Rebase on $I$ and clip at its
length. The extensions agree with the actual maps on $I$, which proves
the assertion without an operand history.
\end{proof}

\begin{definition}[Valid whole-format tagged index]\label{radix:def:valid}
The pieces partition every finite input rank into nonempty intervals and
refine the $2E+1$ original input progressions. Their combined function is
globally nondecreasing and takes values in
$\mathcal F_{\beta,p,E}\cup\{-\infty,+\infty\}$. Finite anchors belong to
$\mathcal F_{\beta,p,E}$, and each staircase height belongs to
$\{1,\beta,\ldots,\beta^{E-1}\}$. At the two infinite inputs, the function
is extended by absorption.
\end{definition}

These requirements hold for the indices constructed from operands below
and are preserved by merging. The local closure lemma permits arbitrary
radix-power heights; the whole-format encoding bound uses the finite
height range in Definition~\ref{radix:def:valid}. Monotonicity within each
piece alone does not imply the global monotonicity required by the sweep.
The bound on the number of new pieces is a general consequence of
monotonicity and interval partitions. The rounding-specific step is closure
in constant-size records whose new parameters can be computed directly.

\begin{corollary}[Whole-format merge and construction]\label{radix:cor:merge}
For the format \eqref{radix:eq:format}, two valid tagged indices with $q_A$ and
$q_B$ pieces can be composed without operands in $O(q_A+q_B)$ abstract
operations, producing at most $\min(N,q_A+q_B-1)$ pieces. For $m\geq1$
operands, balanced construction takes $O(E(1+m\log(m+1)))$ abstract
operations and gives at most $\min(N,4Em+1)$ pieces. A query or inverse
query takes $O(\log(q+1))$ operations. The encoding uses
$O(q(p\log\beta+\log(E+1)))$ bits.
\end{corollary}

\begin{proof}
For a finite affine record the first local index with $f(i)\geq z$ is
$\min(L,\max(0,\lceil(z-c)/d\rceil))$. For a staircase, it is zero
when $z\leq c$. Otherwise set
\[
 k=\left\lceil\frac{z-c}{H}\right\rceil,\qquad
 q=\left\lfloor\frac{k-1}{\lambda}\right\rfloor,\qquad
 s=k-\lambda q.
\]
The inverse is $\min(L,qP+t_s)$. Skipped levels and clipped thresholds
obey the same formula. These constant-operation inverses permit exactly
the sorted-boundary sweep of Theorem~\ref{merge:thm:merge}; replace its
local closure by Lemma~\ref{radix:lem:closure}. At most two inverse
compositions recover the thresholds of each new record.
The composite is globally nondecreasing and format-valued; taking a
maximum of parent heights preserves the finite height range. Thus its
index is valid. A numerical query is located in its original input
progression by binary search and converted to a rank there. This costs
$O(\log(E+1))$, absorbed by $O(\log(q+1))$ since $q\geq2E+1$.

For one addition, add preimages of the $2(E-1)$ interior rounding-region
boundaries and the two overflow cuts to the original $2E+1$ progressions.
This gives at most $4E+1$ pieces in $O(E)$ operations. On each finite
piece rounding has fixed spacing $h$. If $d>h$, or $d=h$ in odd radix,
the translation identity gives an affine record. Otherwise $H=h\geq d$;
monotonicity, the period $\lambda h/d$, and the output lattice give
\eqref{radix:eq:record}. Its thresholds are obtained by solving the
rounding preimage inequalities at $h(k-1/2)$. Their endpoints are
inclusive for even $k$ and strict for odd $k$; in odd radix the endpoints
are unreachable. The exact sums at successive piece starts are increasing;
a single sweep through the ordered rounding regions supplies their
spacings. With that spacing known, quotient and remainder compute the
anchor and each threshold in constant operations. This avoids a fresh
scan of all $E$ regions for each piece. A balanced merge tree gives the stated construction
and piece counts. The stored endpoints, significands, and thresholds
have $O(p\log\beta+\log(E+1))$ bits, and spacing and height exponents
have $O(E)$ possibilities.
\end{proof}

These are abstract-operation bounds. Integer intermediates in least-ulp
units can require $O((p+E)\log\beta)$ bits. No hardware speedup follows.
The records are not required to be unique or bytewise associative.

Both the parity distinction and the operand restriction matter. In the
binary format $p=3,E=2$, the one-operand block $(4)$ maps inputs
$4,5,6,7$ to $8,8,10,12$. A one-threshold height-two record would require
an output increment of two after two input steps, whereas the outputs
at $5$ and $7$ differ by four. For odd radix, allowing arbitrary real
operands also invalidates the one-threshold assertion: uniform
nearest-even rounding to $3\mathbb Z$ gives
\[
 R_3(3i+3/2)=0,6,6,12,12,\ldots\quad(i=0,1,2,3,4,\ldots).
\]
Its increment after one step need not be three. The operand $3/2$ is
not radix-three rational, so this example is excluded by
\eqref{radix:eq:format}.

The initial reference implementation was tested against an enumerated-neighbour
rounding oracle on all operand pairs in seven small formats and all
operand triples in four of them. These checks cover 29,407 pairs,
30,976 triples with both merge associations, and 4,854,320 full-domain
output comparisons. Another 500 blocks include decimal and hexadecimal
formats. Synthetic-record tests check 35,751 compositions and 222,309
outputs without assuming an underlying addition history. These bounded
checks support the proof; they are neither exhaustive checks of standard
formats nor hardware measurements.

A subsequent audit found that the initial singleton implementation
repeated a linear search for the rounding region on each piece. Its
numerical results remain separate from the operation-count claim above.
The corrected implementation uses the single sweep described in the
proof. Independent tests cover 243 blocks in 15 formats, 274,241 forward
comparisons, 4,536 inverse queries, and 165 merges after additional
partition refinement. Another 6,600 synthetic local compositions check
large radix powers and independently shifted record origins. Instrumented
checks exclude the repeated region search from the corrected constructor.
The implementation and the audit commands are listed in the reproducibility
appendix. These tests do not replace the proof of the asymptotic bound.
